\section{The target theory and the embedding discipline}
\label{sec:framework}

\subsection{What is being embedded}

The target is the first-order intuitionistic primitive recursive arithmetic
of \path{ConciseIPRA/Lean/ConciseIPRA/PR.lean} \cite{repo-core}.
Its terms have a symbol for each \emph{specified} primitive recursive
function. These symbols are generated by zero, successor, projections,
composition, and primitive recursion. The nonlogical axioms are the
defining equations of these symbols, successor injectivity,
$Sx\ne0$, and induction for quantifier-free formulas. The logic is
intuitionistic natural deduction with equality, including the usual rules
for $\bot,\wedge,\vee,\to,\forall,\exists$. The quantifier-free induction
schema permits free numerical parameters:
\[
 A(0,\mathbf p)\ \wedge\
 \forall j\bigl(A(j,\mathbf p)\to A(j+1,\mathbf p)\bigr)
 \quad\longrightarrow\quad\forall j\,A(j,\mathbf p).
\]
Here $A$ must be quantifier-free. The theory does not acquire full
induction merely because every function appearing in a formula is PR.

We embed \emph{presentations, observations, and constructions}. A rational
number is an integer numerator and positive denominator; a Gaussian
rational is a pair of rationals; matrices and polynomials are finite
lists. Their equality tests and arithmetic are fixed PR algorithms.
Canonical reduction by gcd is convenient but optional: cross
multiplication gives an exact equality test on valid rational codes.
Malformed inputs receive fixed default outputs so that every numerical
function remains total.

One possible list coding is
\[
 \operatorname{nil}=0,\qquad
 \operatorname{cons}(a,b)=2^a(2b+1).
\]
For a positive code, a bounded search for its exponent of two recovers
the head, and integer division recovers the tail. The tail is smaller
than the code. Iterating for at most the original code therefore
implements length, lookup, and validation by PR functions. Nested lists,
signed integer pairs, and finite tables introduce no new sort into the
first-order theory. This inefficient coding suffices for the logical
claim; the reference implementations use ordinary Python and Lean lists.

An infinite object is described by a fixed PR family of finite
observations. For example, a real name $x(\mathbf p,k)\in\Q$ obeys
\[
 |x(\mathbf p,k)-x(\mathbf p,\ell)|\le2^{-k}+2^{-\ell}.
 \tag{N}\label{eq:regular-name}
\]
The displayed condition is an ordinary first-order formula after coding.
It is not a new real-number sort. Its intended real value is useful
semantic notation; every proof below ultimately concerns finitely many
rational inequalities. Equality of represented objects is generally
unbounded and is never assumed decidable.

There are two different kinds of uniformity. A theorem can be uniform in
the degree, dimension, finite coefficient list, and numerical parameters
of one fixed PR evaluator. Alternatively, it can be a metatheorem
instantiated separately for each PR evaluator. Neither permits one
unclocked PR interpreter for all PR program codes: if such an interpreter
$U$ existed, $n\mapsto U(n,n)+1$ would itself be PR, and its index
would contradict its own evaluation. A bounded universal machine with an
explicit clock is harmless. Our root construction uses fixed coefficient
evaluators and a single rational-polynomial algorithm; the norm
construction has only finite inputs.

\subsection{Why bounded arithmetic is enough for the finite parts}

\begin{lemma}[Bounded formulas have Boolean representatives]
\label{lem:booleanization}
For each fixed bounded numerical formula $B(\mathbf x)$ whose atoms
compare PR numerical terms, one can construct a PR Boolean function
$b_B$ such that
\[
 \ipra\ \vdash\ B(\mathbf x)\ \longleftrightarrow\ b_B(\mathbf x)=1.
\]
This is a metatheorem for fixed syntax, not an evaluator for arbitrary
formulas and arbitrary PR symbol codes.
\end{lemma}

\begin{proof}
Natural-number equality and order have PR characteristic functions;
their specifications follow by quantifier-free induction on the usual
recursive comparison functions. Rational comparisons reduce to integer
comparisons after multiplication by positive denominators. Boolean
negation, conjunction, and disjunction handle the corresponding
connectives. Since the already constructed Boolean atoms are decidable,
implication is represented by $\neg b\vee c$: the verification is
intuitionistic case analysis on a Boolean value, not excluded middle for
arbitrary propositions.

For $\exists y<t\,C(\mathbf x,y)$, take the finite Boolean disjunction
of $b_C(\mathbf x,0),\ldots,b_C(\mathbf x,t-1)$. For a bounded universal
quantifier, take the finite conjunction. The empty disjunction is false
and the empty conjunction true. These equivalences can be verified
without assuming bounded induction in advance. Write $A(i)$ for the
conjunction of the first $i$ Boolean values. From the bounded-universal
premise for a fixed terminal $T$, use the quantifier-free guarded
invariant $i\le T\to A(i)=1$. Conversely, fix an index $j$ and induct
on $i$ with the quantifier-free invariant
$j<i\wedge A(i)=1\to b_C(\mathbf x,j)=1$; split the successor case
into $j=i$ and $j<i$ by numerical decidability. After this induction,
introduce the bounded universal quantifier over $j$.

For the existential fold $E(i)$, maintain a first-success register
$w(i)$, with default $i$ when there has been no success. Its recurrence
either retains the old successful index or records $i$ on the next
successful test. The invariants
\[
 E(i)=1\ \to\ w(i)<i\wedge b_C(\mathbf x,w(i))=1,
 \qquad
 j<i\wedge b_C(\mathbf x,j)=1\ \to\ E(i)=1
\]
are quantifier-free, with $j$ a fixed parameter in the second induction.
They prove both implications and produce the existential witness.
Each fixed nesting of bounded quantifiers adds a fixed nesting of PR
recursions. Structural induction on the fixed formula completes the
construction. The detailed arithmetic groundwork is also provided by
the earlier finite-contract treatment \cite{repo-contracts}.
\end{proof}

In particular, induction on a bounded formula is an admissible derived
rule: replace it by its Boolean equation, use quantifier-free induction,
and translate back. To reason about a loop with terminal length $T$, use
the guarded invariant
\[
                         j\le T\ \longrightarrow\ I(j).
\]
Its Boolean representative is quantifier-free. For $j+1\le T$ the
local transition establishes the next invariant; when this guard fails,
the implication holds by decidable arithmetic. Thus no induction beyond
the terminal index is required.

The surrounding hypothesis context may contain unbounded validity
assertions, for example \eqref{eq:regular-name}. Such assertions remain
hypotheses. At a finite step one instantiates them at the finitely many
needed indices; one does not put them into the induction invariant.
If all those indices have a PR bound, their needed instances can instead
be collected into one bounded Boolean table. This is the precise role of
the observation-demand calculations that follow.

\subsection{Five separate obligations}

The previous frameworks already established bounded Booleanization,
guarded loops, certificate composition, finite observation contracts,
and several quantitative representation tools
\cite{repo-contracts,repo-embedding,repo-first-framework}. We retain that
foundation and make the following obligations explicit for each new
construction.
\begin{enumerate}
\item \textbf{Representation.} Specify the finite codes, validity
conditions, and what equality of the represented objects means.
\item \textbf{Production.} Give a fixed PR algorithm, including bounds
for every search and a default on invalid inputs. A PR checker alone
does not supply a passing certificate.
\item \textbf{Finite correctness.} Prove the output observations from
identified input observations. Supply PR demand bounds and bounded
induction invariants.
\item \textbf{Adequacy.} Explain exactly which mathematical conclusion
the observations express. Small residuals, a coherent witness, a finite
set of witnesses, and a complete list with a height bound are different
outputs.
\item \textbf{Formalization.} Construct the object-theory derivation,
or label the remaining derivation as a target. Native Lean correctness
and an embedded \texttt{Prf PRAAx} term are distinct achievements.
\end{enumerate}

There is no general rule that a true universal statement about PR
functions is provable in \ipra. The displayed proofs supply the needed
finite derivations; their eventual Lean compilation must reproduce
those derivations. Likewise, a complete analytic statement may still
require a new estimate. The framework organizes such an estimate and
its uses; it does not manufacture it from the name of a certificate.

\subsection{How the new components fit together}

Section~\ref{sec:backward-demands} handles a family of constructions whose
correctness at precision $k$ depends on an earlier stage at a different
precision. The proof fixes the final demand and computes all intermediate
demands backward, so its induction invariant remains bounded.

Section~\ref{sec:aggregates} supplies a finite-set interface when the task
requires coverage of possible witnesses but no canonical choice of one.
It gives arithmetic proofs of composition, distances, images, and
approximate optimization. Section~\ref{sec:rootsets} supplies a producer
for that interface from monic polynomial coefficient names, including
multiple roots.

Section~\ref{sec:norm-certificates} addresses a different obstacle:
turning a finite algebraic presentation into a certificate with enough
mathematical force to bound a fiber. A determinant can be constructed
uniformly, and its nonvanishing can be proved automatically for coprime
superelliptic exponents. This supplies a previously missing geometric
input to a Diophantine proof strategy; it does not supply that strategy's
remaining approximation theorem or global height bound.
